% =====================================================================
% Topic 5 — Logic and Puzzles — part c: problems 31–40 (extension)
% =====================================================================
% ---- helper macros for parts c/d of topic 5 (prefix \lgxcd, guarded) ----
% shape codes: 1 = circle, 2 = square, 3 = triangle, 4 = star
\providecommand{\lgxcdShape}[3]{\ifcase#1\or
  \draw[Ink,fill=Slate!35,line width=0.6pt] (#2,#3) circle (0.19);\or
  \draw[Ink,fill=tmBlue!35,line width=0.6pt] ($(#2,#3)+(-0.18,-0.18)$) rectangle ($(#2,#3)+(0.18,0.18)$);\or
  \draw[Ink,fill=tmOrange!45,line width=0.6pt] ($(#2,#3)+(-0.2,0.17)$)--($(#2,#3)+(0.2,0.17)$)--($(#2,#3)+(0,-0.19)$)--cycle;\or
  \draw[Ink,fill=tmGreen!45,line width=0.6pt,shift={(#2,#3)}] (90:0.24) \foreach \lgxcdk in {1,...,9}{ -- ({90+36*\lgxcdk}:{ifthenelse(mod(\lgxcdk,2)==1,0.1,0.24)})} -- cycle;\fi}
% inline shape in text
\providecommand{\lgxcdI}[1]{\tikz[baseline=-0.7ex]{\lgxcdShape{#1}{0}{0}}}
% a row of n shapes centred at the origin: \lgxcdRow{n}{codes}
\providecommand{\lgxcdRow}[2]{\foreach \lgxcdt [count=\lgxcdi] in {#2}{%
  \pgfmathsetmacro\lgxcdxx{(\lgxcdi-(#1+1)/2)*0.48}\lgxcdShape{\lgxcdt}{\lgxcdxx}{0.21}}}
% equal-arm balance at height #1: \lgxcdBal{y}{n left}{codes left}{n right}{codes right}
\providecommand{\lgxcdBal}[5]{\begin{scope}[yshift=#1cm]
  \draw[Ink,line width=1pt] (-2.7,0)--(-0.3,0) (0.3,0)--(2.7,0);
  \draw[Ink,line width=0.9pt] (-1.5,0)--(-1.5,-0.55)--(1.5,-0.55)--(1.5,0);
  \draw[Ink,line width=0.9pt] (-0.18,-0.85)--(0,-0.58)--(0.18,-0.85);
  \begin{scope}[xshift=-1.5cm]\lgxcdRow{#2}{#3}\end{scope}
  \begin{scope}[xshift=1.5cm]\lgxcdRow{#4}{#5}\end{scope}\end{scope}}
% square grid n x n with box lines every #2 rows / #3 columns; #4 = given r/c/v, #5 = filled r/c/v (rows from top)
\providecommand{\lgxcdSud}[5]{\begin{scope}[x=0.62cm,y=0.62cm]
  \foreach \lgxcdi in {0,...,#1}{\draw[Slate!60,line width=0.5pt] (0,\lgxcdi)--(#1,\lgxcdi) (\lgxcdi,0)--(\lgxcdi,#1);}
  \foreach \lgxcdi in {0,#2,...,#1}{\draw[Ink,line width=1.4pt] (0,\lgxcdi)--(#1,\lgxcdi);}
  \foreach \lgxcdi in {0,#3,...,#1}{\draw[Ink,line width=1.4pt] (\lgxcdi,0)--(\lgxcdi,#1);}
  \foreach \lgxcdr/\lgxcdc/\lgxcdv in {#4}{\node[font=\sffamily\bfseries,text=Ink] at (\lgxcdc+0.5,#1-\lgxcdr-0.5){\lgxcdv};}
  \foreach \lgxcdr/\lgxcdc/\lgxcdv in {#5}{\node[font=\sffamily,text=tmBlue] at (\lgxcdc+0.5,#1-\lgxcdr-0.5){\lgxcdv};}
  \end{scope}}

\clearpage
\dividerbar{Extension and challenge}

% ---------------------------------------------------------------- Problem 31
\begin{probbox}{31}{\relax}
On an island, every inhabitant is either a \emph{truth-teller} (always tells the truth) or a \emph{liar} (always lies).
\begin{enumerate}[label=\alph*),leftmargin=1.6em,itemsep=1pt,topsep=2pt]
  \item You meet two inhabitants, Khánh and Lâm. Khánh says: ``Both of us are liars.'' What type is each of them?
  \item You meet three inhabitants, An, Bình and Chi.\\
  An says: ``Bình is a liar.''\\
  Bình says: ``Chi is a liar.''\\
  Chi says: ``An and Bình are both liars.''\\
  What type is each of them?
\end{enumerate}
\end{probbox}
\probsol{
\textbf{Understanding the problem.} Each person has only two possibilities. The safest method is to \emph{assume} one possibility and
check whether the statements fit that assumption. If we reach a contradiction, we rule that possibility out.
\begin{recallbox}Whatever a truth-teller says is \textbf{true}. Whatever a liar says is \textbf{false}. So each statement gives us an
``equation'': \emph{the statement is true} $\Longleftrightarrow$ \emph{the speaker is a truth-teller}.\end{recallbox}
\par\smallskip\textbf{Part a).}
\step{1}{Suppose Khánh is a truth-teller}
Then his statement is true, so Khánh is a liar -- contradicting the assumption. So Khánh \textbf{cannot} be a truth-teller.
\step{2}{Khánh is a liar}
The statement ``Both of us are liars'' is false. It is false when \emph{at least one} of them is not a liar. Khánh is a liar, so the
other one -- Lâm -- must be a truth-teller.
\begin{warnbox}The opposite of ``both are liars'' is \textbf{not} ``both tell the truth'', but ``at least one of them tells the
truth''.\end{warnbox}
\par\smallskip\textbf{Part b).}
\step{3}{Suppose An is a truth-teller}
An's statement is true $\Rightarrow$ Bình is a liar $\Rightarrow$ Bình's statement is false $\Rightarrow$ Chi is a truth-teller
$\Rightarrow$ Chi's statement is true $\Rightarrow$ An is a liar. This contradicts the assumption! Ruled out.
\step{4}{So An is a liar}
An's statement is false $\Rightarrow$ Bình is \emph{not} a liar, i.e.\ Bình is a truth-teller $\Rightarrow$ Bình's statement is true
$\Rightarrow$ Chi is a liar.
\step{5}{Check Chi's statement}
Chi is a liar, so the statement ``An and Bình are both liars'' must be false. Indeed it is false, because Bình is a truth-teller.
There is no contradiction.
\begin{center}
\begin{tabular}{|c|c|l|c|}
\hline
\rowcolor{tmBlue!10} Person & Type & Statement & True/false\\ \hline
An & liar & Bình is a liar & false \\ \hline
Bình & truth-teller & Chi is a liar & true\\ \hline
Chi & liar & An and Bình are both liars & false\\ \hline
\end{tabular}
\end{center}
\verify{In the table, every truth-teller makes a true statement and every liar a false one. We have also seen that ``An is a
truth-teller'' is impossible, so this is the only solution.}
\ans{a) Khánh is a liar, Lâm a truth-teller;\quad b) An and Chi are liars, Bình is a truth-teller}
}

% ---------------------------------------------------------------- Problem 32
\begin{probbox}{32}{\relax}
Mai lies on Mondays, Tuesdays and Wednesdays and tells the truth on the other days of the week. Hùng lies on Thursdays, Fridays and
Saturdays and tells the truth on the other days. One day, both of them say: ``Yesterday was a day on which I lie.'' Which day of the
week is it?
\end{probbox}
\probsol{
\textbf{Understanding the problem.} The statement ``Yesterday was a day on which I lie'' may be true or false depending on the day.
For each person we find the days on which they \emph{could} say it, then take the day common to both.
\begin{recallbox}On a truth day the statement must be true; on a lying day the statement must be false.\end{recallbox}
\step{1}{Table for Mai (lies Mon, Tue, Wed)}
\begin{center}\small
\begin{tabular}{|l|c|c|c|c|c|c|c|}
\hline
\rowcolor{tmBlue!10} Today & Mon & Tue & Wed & Thu & Fri & Sat & Sun\\ \hline
Mai today & lies & lies & lies & truth & truth & truth & truth\\ \hline
Mai yesterday & truth (Sun) & lies & lies & lies & truth & truth & truth\\ \hline
Statement true? & no & yes & yes & yes & no & no & no\\ \hline
Can Mai say it? & \textbf{yes} & no & no & \textbf{yes} & no & no & no\\ \hline
\end{tabular}
\end{center}
Mai can say it when she lies today and the statement is false (Monday), or when she tells the truth today and the statement is true
(Thursday).
\step{2}{Table for Hùng (lies Thu, Fri, Sat)}
\begin{center}\small
\begin{tabular}{|l|c|c|c|c|c|c|c|}
\hline
\rowcolor{tmBlue!10} Today & Mon & Tue & Wed & Thu & Fri & Sat & Sun\\ \hline
Hùng today & truth & truth & truth & lies & lies & lies & truth\\ \hline
Hùng yesterday & truth & truth & truth & truth & lies & lies & lies\\ \hline
Statement true? & no & no & no & no & yes & yes & yes\\ \hline
Can Hùng say it? & no & no & no & \textbf{yes} & no & no & \textbf{yes}\\ \hline
\end{tabular}
\end{center}
Hùng can say it on Thursday or Sunday.
\step{3}{The common day}
Mai: $\{$Monday, Thursday$\}$; Hùng: $\{$Thursday, Sunday$\}$. The only common day is \textbf{Thursday}.
\begin{tipbox}Quick rule: a person can say ``Yesterday I lied'' only on a \emph{switch day} (yesterday and today are of different
types): the first lying day, or the first day back to telling the truth.\end{tipbox}
\verify{Thursday: Mai tells the truth, and yesterday (Wednesday) she lied -- the statement is true, consistent. Hùng lies, and
yesterday (Wednesday) he told the truth -- the statement is false, consistent with Hùng lying today.}
\ans{Thursday}
}

% ---------------------------------------------------------------- Problem 33
\begin{probbox}{33}{\relax}
Six friends Lan, Minh, Nam, Oanh, Phúc and Quân run a race. No two of them finish at the same time. We know that:
\begin{itemize}[leftmargin=1.4em,itemsep=0pt,topsep=2pt]
  \item Minh finishes after Lan but before Nam;
  \item Oanh finishes after Nam;
  \item Phúc finishes \emph{immediately after} Lan;
  \item Quân is neither first nor last, and Quân finishes before Minh.
\end{itemize}
Find the finishing order of the six friends. Who finishes third?
\end{probbox}
\probsol{
\textbf{Understanding the problem.} We have ``before -- after'' conditions. A good method is to join them into a \emph{chain} and then
insert the remaining friends one by one. Write $X \to Y$ for ``$X$ finishes before $Y$''.
\step{1}{Join conditions 1 and 2}
Lan $\to$ Minh $\to$ Nam and Nam $\to$ Oanh, so: \ Lan $\to$ Minh $\to$ Nam $\to$ Oanh.
\step{2}{Insert Phúc}
Phúc finishes \emph{immediately after} Lan, i.e.\ nobody comes between Lan and Phúc: \ Lan, Phúc $\to$ Minh $\to$ Nam $\to$ Oanh.
\step{3}{Insert Quân}
Quân finishes before Minh, so Quân is in one of these places: before Lan, or between Phúc and Minh (he may not come between Lan and
Phúc). But Quân is not first, so he is not before Lan. Hence:
\[ \text{Lan} \to \text{Phúc} \to \text{Quân} \to \text{Minh} \to \text{Nam} \to \text{Oanh}. \]
\begin{figC}
\foreach \n/\x in {Lan/0,Phúc/1,Quân/2,Minh/3,Nam/4,Oanh/5}{
  \node[draw=tmBlue,fill=tmBlue!8,rounded corners=3pt,minimum width=1.5cm,minimum height=0.7cm] at (\x*1.9,0){\n};}
\foreach \x [count=\k] in {0,...,5}{\node[font=\scriptsize,text=Slate] at (\x*1.9,-0.6){place \k};}
\foreach \x in {0,...,4}{\draw[->,tmBlue,thick] (\x*1.9+0.78,0)--(\x*1.9+1.12,0);}
\end{figC}
\step{4}{Check every condition again}
Minh (4th) is after Lan (1st) and before Nam (5th): true. Oanh (6th) is after Nam: true. Phúc (2nd) is immediately after Lan: true.
Quân is 3rd -- not first, not last, before Minh: true.
\begin{warnbox}If you forget the word ``immediately'', Quân could come between Lan and Phúc and the problem would have two answers. Read
every word of each condition carefully!\end{warnbox}
\verify{We tried every possible position for Quân; only one order satisfies all four conditions.}
\ans{Lan, Phúc, Quân, Minh, Nam, Oanh; Quân finishes third}
}

% ---------------------------------------------------------------- Problem 34
\begin{probbox}{34}{\relax}
Five friends Hà, Khoa, Linh, My and Tùng sit in a row of $5$ chairs numbered $1, 2, 3, 4, 5$ from left to right. We know that:
Hà and Khoa sit next to each other; Linh and My do \emph{not} sit next to each other; Tùng sits at one end of the row (chair $1$ or
chair $5$).
\begin{enumerate}[label=\alph*),leftmargin=1.6em,itemsep=1pt,topsep=2pt]
  \item How many seating arrangements satisfy these conditions?
  \item Which of the following statements is \emph{certainly true}?\\
  (1) Linh sits at one end of the row. \quad (2) Hà does not sit at an end of the row. \quad (3) Khoa sits in chair $3$.
  \item If we also know that Linh sits in chair $1$ and Hà sits to the left of Khoa, find each friend's seat.
\end{enumerate}
\end{probbox}
\probsol{
\textbf{Understanding the problem.} Start with the ``tightest'' condition (Tùng has only $2$ places), then look at the $4$ remaining
chairs.
\step{1}{Case: Tùng sits in chair 5}
The other four sit in chairs $1, 2, 3, 4$. The pair Hà -- Khoa takes two adjacent chairs: $(1,2)$, $(2,3)$ or $(3,4)$.
\begin{itemize}[leftmargin=1.4em,itemsep=0pt,topsep=1pt]
  \item Hà -- Khoa in chairs $(1,2)$: Linh and My must sit in chairs $3, 4$ -- next to each other, \textbf{impossible}.
  \item Hà -- Khoa in chairs $(3,4)$: Linh and My sit in chairs $1, 2$ -- next to each other, \textbf{impossible}.
  \item Hà -- Khoa in chairs $(2,3)$: Linh and My sit in chairs $1$ and $4$ -- not next to each other, \textbf{possible}.
\end{itemize}
In the possible case, Hà and Khoa can swap ($2$ ways) and Linh and My can swap ($2$ ways): $2\times 2 = 4$ ways.
\step{2}{Case: Tùng sits in chair 1}
Exactly the same (turn the row around): Hà -- Khoa in chairs $(3,4)$, Linh and My in chairs $2$ and $5$: also $4$ ways.
\step{3}{Total number of ways}
$4 + 4 = 8$ ways.
\begin{recallbox}\textbf{Multiplication rule:} if a task consists of consecutive steps, step 1 in $m$ ways and step 2 in $n$ ways, then
the task can be done in $m\times n$ ways. \textbf{Addition rule:} separate cases are added.\end{recallbox}
\step{4}{Part b): check each statement}
\begin{center}\small
\begin{tabular}{|c|c|c|c|c|}
\hline
\rowcolor{tmBlue!10} Chair 1 & Chair 2 & Chair 3 & Chair 4 & Chair 5\\ \hline
Linh/My & Hà/Khoa & Hà/Khoa & My/Linh & Tùng\\ \hline
Tùng & Linh/My & Hà/Khoa & Hà/Khoa & My/Linh\\ \hline
\end{tabular}
\end{center}
(1) False: when Tùng is in chair 1, Linh may sit in chair 2. \ (2) \textbf{True}: in both rows, Hà and Khoa are always in middle chairs
($2, 3$ or $3, 4$). \ (3) False: Khoa may sit in chair $2$ or chair $4$.
\step{5}{Part c)}
Linh is in chair $1$, so Tùng is not in chair $1$; Tùng is in chair $5$ (first row of the table). Then My is in chair $4$ and Hà --
Khoa in chairs $2, 3$; Hà is to the left of Khoa, so Hà is in chair $2$ and Khoa in chair $3$.
\verify{Linh(1), Hà(2), Khoa(3), My(4), Tùng(5): Hà and Khoa are adjacent; Linh in chair 1 and My in chair 4 are not adjacent; Tùng
is at an end. Correct.}
\ans{a) $8$ ways;\quad b) statement (2);\quad c) Linh -- Hà -- Khoa -- My -- Tùng (chairs 1 to 5)}
}

% ---------------------------------------------------------------- Problem 35
\begin{probbox}{35}{\relax}
Number the letters $A, B, C, \dots, Z$ (the $26$-letter English alphabet) as $0, 1, 2, \dots, 25$. The \emph{Caesar cipher with key
$k$} replaces each letter with number $x$ by the letter whose number is the \emph{remainder} of $x + k$ when divided by $26$
(spaces are kept).
\begin{figC}
\foreach \L [count=\i from 0] in {A,B,C,D,E,F,G,H,I,J,K,L,M,N,O,P,Q,R,S,T,U,V,W,X,Y,Z}{
  \node[font=\sffamily\small\bfseries,text=Ink] at ({90-\i*360/26}:2.15){\L};
  \node[font=\sffamily\tiny,text=tmBlue] at ({90-\i*360/26}:1.65){\i};}
\draw[Line,line width=0.8pt] (0,0) circle (1.9);
\draw[->,tmOrange,line width=1pt] (110:1.2) arc (110:30:1.2);
\node[font=\scriptsize,text=tmOrange!80!black,align=center] at (0,0){move\\clockwise};
\end{figC}
\begin{enumerate}[label=\alph*),leftmargin=1.6em,itemsep=1pt,topsep=2pt]
  \item Encrypt the text \textsf{TOAN HOC} with key $k = 5$.
  \item Decrypt the text \textsf{OVJ CBP}, knowing that the key is $k = 7$.
  \item In a message, the word \textsf{BAN} was encrypted as \textsf{QPC}. Find the key $k$ (with $0 \le k \le 25$), then decrypt the
  word \textsf{HPJ} from the same message.
  \item A message is encrypted with key $10$, and the result is encrypted once more with key $20$. These two encryptions together
  are the same as \emph{one} encryption with which key?
  \item Tú encrypts a word with key $4$, then encrypts the result with key $4$ again, and so on. After how many encryptions, at least,
  does Tú get back exactly the original word?
\end{enumerate}
\end{probbox}
\probsol{
\textbf{Understanding the problem.} When adding the key takes us past $25$, we ``wrap around'' to the start of the alphabet -- that is
exactly taking the \emph{remainder on division by $26$}. Decrypting is the reverse: \emph{subtract} the key (if the result is
negative, add $26$).
\begin{recallbox}The remainder of $a$ on division by $26$ is the number $r$ with $0 \le r \le 25$ such that $a - r$ is divisible by
$26$. Examples: $30$ leaves remainder $4$; \ $-6$ ``leaves remainder'' $20$, because $-6 + 26 = 20$.\end{recallbox}
\step{1}{Part a): encrypt \textsf{TOAN HOC} with $k = 5$}
\begin{center}\small
\begin{tabular}{|c|c|c|c|c|c|c|c|}
\hline
\rowcolor{tmBlue!10} Plain letter & T & O & A & N & H & O & C\\ \hline
Number $x$ & $19$ & $14$ & $0$ & $13$ & $7$ & $14$ & $2$\\ \hline
$x + 5$ & $24$ & $19$ & $5$ & $18$ & $12$ & $19$ & $7$\\ \hline
Cipher letter & Y & T & F & S & M & T & H\\ \hline
\end{tabular}
\end{center}
Result: \textsf{YTFS MTH}. (No number exceeds $25$ here, so no wrapping is needed yet.)
\step{2}{Part b): decrypt \textsf{OVJ CBP} with $k = 7$ (subtract $7$)}
\begin{center}\small
\begin{tabular}{|c|c|c|c|c|c|c|}
\hline
\rowcolor{tmBlue!10} Cipher letter & O & V & J & C & B & P\\ \hline
Number & $14$ & $21$ & $9$ & $2$ & $1$ & $15$\\ \hline
Number $- 7$ & $7$ & $14$ & $2$ & $-5 \to 21$ & $-6 \to 20$ & $8$\\ \hline
Plain letter & H & O & C & V & U & I\\ \hline
\end{tabular}
\end{center}
Original text: \textsf{HOC VUI} (Vietnamese for ``fun learning''). Note: $2 - 7 = -5$, adding $26$ gives $21$ (the letter V).
\step{3}{Part c): find the key from the pair \textsf{BAN} $\to$ \textsf{QPC}}
B ($1$) becomes Q ($16$), so $k = 16 - 1 = 15$. Check: A ($0$) $\to 0 + 15 = 15$ (P), correct; N ($13$) $\to 13 + 15 = 28$, remainder
$2$ on division by $26$ (C), correct. So $k = 15$.\\
Decrypt \textsf{HPJ} (subtract $15$): H $= 7$, $7 - 15 = -8 \to 18$ (S); \ P $= 15 \to 0$ (A); \ J $= 9$, $9 - 15 = -6 \to 20$ (U).
The original word is \textsf{SAU}.
\step{4}{Part d): two encryptions in a row}
Each letter gets $10$ added and then $20$ added, i.e.\ $30$ added. Since $30$ leaves remainder $4$ on division by $26$, ``add $30$ and
take the remainder'' is exactly the same as ``add $4$ and take the remainder''. The two encryptions equal \textbf{one encryption with
key $4$}.
\step{5}{Part e): repeated encryption with key $4$}
After $n$ encryptions, every letter has moved $4n$ places. To get back to the original letter we must go round a whole number of
times, i.e.\ $4n$ must be divisible by $26$. Try $n = 1, 2, 3, \dots$: \ $4n = 4, 8, 12, \dots$ The smallest (non-zero) $4n$
divisible by $26$ is $\text{LCM}(4, 26) = 52$, which gives $n = 52 \div 4 = 13$.
\begin{recallbox}$4n$ is divisible by $26 = 2\times 13$ $\Leftrightarrow$ $2n$ is divisible by $13$ $\Leftrightarrow$ $n$ is divisible
by $13$ (because $13$ is a prime that does not divide $2$).\end{recallbox}
\begin{tipbox}With key $13$ you only need $2$ encryptions to get back ($2\times 13 = 26$). Encrypting twice with key $13$ is the same as
decrypting!\end{tipbox}
\verify{Part a): decrypting \textsf{YTFS MTH} by subtracting $5$ gives back \textsf{TOAN HOC}. Part c): \textsf{SAU} encrypted with key
$15$: S $18 + 15 = 33 \to 7$ (H), A $\to 15$ (P), U $20 + 15 = 35 \to 9$ (J). Indeed \textsf{HPJ}.}
\ans{a) \textsf{YTFS MTH};\ b) \textsf{HOC VUI};\ c) $k = 15$, \textsf{SAU};\ d) key $4$;\ e) $13$ times}
}

% ---------------------------------------------------------------- Problem 36
\begin{probbox}{36}{\stars{1}}
Again number $A = 0, B = 1, \dots, Z = 25$. A \emph{multiplicative cipher} replaces the letter with number $x$ by the letter whose number
is the remainder of $3x + 2$ on division by $26$.
\begin{enumerate}[label=\alph*),leftmargin=1.6em,itemsep=1pt,topsep=2pt]
  \item Encrypt the word \textsf{MEO}.
  \item Decrypt the word \textsf{ECS}.
  \item Nhật suggests using $2x + 2$ instead of $3x + 2$. Find two different letters with the same cipher letter. Why does that make
  decryption impossible?
  \item Explain why, with $3x + 2$, two different letters always get two different cipher letters.
\end{enumerate}
\end{probbox}
\probsol{
\textbf{Understanding the problem.} This cipher does not just ``shift'': it also multiplies and then takes the remainder. Encrypting is
easy (compute directly); to decrypt we need $x$ when we know the remainder of $3x + 2$ -- we use the trick ``add $26$ until the number
is divisible by $3$''.
\step{1}{Part a): encrypt \textsf{MEO}}
M $= 12$: $3\cdot 12 + 2 = 38$, remainder $12$ on division by $26$ $\to$ M.\\
E $= 4$: $3\cdot 4 + 2 = 14$ $\to$ O.\\
O $= 14$: $3\cdot 14 + 2 = 44$, remainder $18$ $\to$ S.\\
Result: \textsf{MOS}. (Fun fact: M is encrypted as itself.)
\step{2}{Part b): decrypt \textsf{ECS}}
For a cipher letter with number $y$ we need $x$ ($0 \le x \le 25$) such that $3x + 2$ leaves remainder $y$, i.e.\ $3x = y - 2$, or
$y - 2 + 26$, or $y - 2 + 52$ (try them in turn and take the one divisible by $3$).
\begin{center}\small
\begin{tabular}{|c|c|c|c|c|}
\hline
\rowcolor{tmBlue!10} Cipher letter & $y$ & Try $y - 2,\ y + 24,\ y + 50$ & $3x$ & Plain letter\\ \hline
E & $4$ & $2$ (no), $28$ (no), $54$ (divisible) & $54 \Rightarrow x = 18$ & S\\ \hline
C & $2$ & $0$ (divisible) & $0 \Rightarrow x = 0$ & A\\ \hline
S & $18$ & $16$ (no), $42$ (divisible) & $42 \Rightarrow x = 14$ & O\\ \hline
\end{tabular}
\end{center}
The original word is \textsf{SAO}.
\step{3}{Part c): why $2x + 2$ does not work}
A $= 0$: $2\cdot 0 + 2 = 2$ $\to$ C. \quad N $= 13$: $2\cdot 13 + 2 = 28$, remainder $2$ $\to$ C.\\
Both A and N become C. When the receiver sees a C, they cannot tell whether the original letter was A or N, so decryption is not
certain. (In fact every number $2x + 2$ is even, so there are only $13$ cipher letters for $26$ plain letters -- some pairs must
coincide.)
\step{4}{Part d): why $3x + 2$ never collides}
\begin{proofbox}Suppose two letters $x$ and $z$ (with $0 \le z < x \le 25$) have the same cipher letter. Then $3x + 2$ and $3z + 2$ leave
the same remainder on division by $26$, so their difference $3(x - z)$ is divisible by $26$. But $3$ and $26$ have no common divisor
greater than $1$, so $x - z$ must be divisible by $26$. However $0 < x - z \le 25$ -- not divisible by $26$. Contradiction. So
different letters always get different cipher letters.\end{proofbox}
\begin{warnbox}With $2x + 2$, the argument breaks down here: $2(x - z)$ is divisible by $26$ as soon as $x - z$ is divisible by $13$,
for example $x - z = 13$ (A and N).\end{warnbox}
\verify{Encrypt \textsf{SAO} again: S $3\cdot 18 + 2 = 56 \to 56 - 52 = 4$ (E); A $\to 2$ (C); O $\to 44 - 26 = 18$ (S). We get
\textsf{ECS}, correct.}
\ans{a) \textsf{MOS};\ b) \textsf{SAO};\ c) A and N both become C;\ d) see the proof}
}

% ---------------------------------------------------------------- Problem 37
\begin{probbox}{37}{\relax}
In a $3\times 3$ \emph{magic square}, the numbers in every row, every column and both diagonals have the same sum (called the
\emph{magic sum} $M$).
\begin{enumerate}[label=\alph*),leftmargin=1.6em,itemsep=1pt,topsep=2pt]
  \item Let $c$ be the number in the centre cell. Prove that $M = 3c$.
  \item Fill in the empty cells of the magic square below.
\end{enumerate}
\begin{figC}
\lgxcdSud{3}{3}{3}{0/0/9,0/1/19,2/2/13}{}
\end{figC}
\end{probbox}
\probsol{
\textbf{Understanding the problem.} In part b) only three numbers are known; we know neither the centre nor $M$. Part a) is the ``key'':
it links the centre number to $M$.
\step{1}{Part a): add the four lines through the centre}
Name the cells:
\begin{figC}
\lgxcdSud{3}{3}{3}{0/0/$a$,0/1/$b$,0/2/$d$,1/0/$e$,1/1/$c$,1/2/$f$,2/0/$g$,2/1/$h$,2/2/$i$}{}
\end{figC}
Four lines pass through the centre: the middle row $e + c + f$, the middle column $b + c + h$, and the two diagonals $a + c + i$ and
$d + c + g$. Adding these four sums:
\[ (e + c + f) + (b + c + h) + (a + c + i) + (d + c + g) = 4M. \]
On the left, each outer cell appears exactly \emph{once}, while the centre $c$ appears \emph{four} times. So the left side equals
(sum of all $9$ cells) $+\, 3c$.
\begin{proofbox}The sum of all $9$ cells is the sum of the three rows, i.e.\ $3M$. Therefore $3M + 3c = 4M$, which gives $3c = M$.\end{proofbox}
\step{2}{Part b): find the centre}
The main diagonal passes through $9$, the centre $c$ and $13$: \ $9 + c + 13 = M = 3c$. So $22 = 2c$, $c = 11$ and $M = 33$.
\step{3}{Fill the cells one by one (each time choose a line with only one cell missing)}
\begin{itemize}[leftmargin=1.4em,itemsep=0pt,topsep=1pt]
  \item Top row: $33 - 9 - 19 = 5$ (top right corner).
  \item Other diagonal: $33 - 5 - 11 = 17$ (bottom left corner).
  \item Left column: $33 - 9 - 17 = 7$; \quad middle row: $33 - 7 - 11 = 15$.
  \item Bottom row: $33 - 17 - 13 = 3$.
\end{itemize}
\begin{figC}
\lgxcdSud{3}{3}{3}{0/0/9,0/1/19,2/2/13}{0/2/5,1/0/7,1/1/11,1/2/15,2/0/17,2/1/3}
\end{figC}
\verify{Middle column $19 + 11 + 3 = 33$; right column $5 + 15 + 13 = 33$; middle row $7 + 11 + 15 = 33$; the diagonals $9 + 11 + 13$ and
$5 + 11 + 17$ both equal $33$.}
\begin{tipbox}The result $M = 3c$ holds for \emph{every} $3\times 3$ magic square, even with negative numbers or fractions.\end{tipbox}
\ans{a) see the proof;\quad b) top row $9, 19, 5$; middle row $7, 11, 15$; bottom row $17, 3, 13$}
}

% ---------------------------------------------------------------- Problem 38
\begin{probbox}{38}{\relax}
The three equal-arm balances below are all balanced. Objects of the same shape have the same mass.
\begin{figC}
\lgxcdBal{0}{2}{2,3}{5}{1,1,1,1,1}
\lgxcdBal{-1.5}{2}{3,3}{4}{1,1,1,1}
\lgxcdBal{-3}{2}{4,1}{2}{2,2}
\node[font=\small\bfseries,text=Slate] at (-3.3,0){(1)};
\node[font=\small\bfseries,text=Slate] at (-3.3,-1.5){(2)};
\node[font=\small\bfseries,text=Slate] at (-3.3,-3){(3)};
\end{figC}
\begin{enumerate}[label=\alph*),leftmargin=1.6em,itemsep=1pt,topsep=2pt]
  \item How many \lgxcdI{1} does each of \lgxcdI{3}, \lgxcdI{2}, \lgxcdI{4} weigh?
  \item How many \lgxcdI{3} are needed to balance \lgxcdI{4}\,\lgxcdI{4}\,\lgxcdI{2}\,\lgxcdI{2}?
  \item Each \lgxcdI{1} weighs $12$ g. How many grams do the three objects \lgxcdI{4}\,\lgxcdI{3}\,\lgxcdI{2} weigh altogether?
\end{enumerate}
\end{probbox}
\probsol{
\textbf{Understanding the problem.} Choose \lgxcdI{1} as the ``unit of measure''. Look for a balance with only \emph{one kind} of unknown
object, solve it first, and substitute into the other balances.
\begin{recallbox}A balanced scale is like an equation. We may remove (or add) identical objects on both pans, or split both pans into
equal parts, and the scale stays balanced.\end{recallbox}
\step{1}{Balance (2): find \lgxcdI{3}}
Two \lgxcdI{3} balance four \lgxcdI{1}. Halve each pan: one \lgxcdI{3} balances two \lgxcdI{1}. \quad \lgxcdI{3} $= 2$ \lgxcdI{1}.
\step{2}{Balance (1): find \lgxcdI{2}}
\lgxcdI{2}\,\lgxcdI{3} balance five \lgxcdI{1}. Replace \lgxcdI{3} by two \lgxcdI{1}: \lgxcdI{2} together with two \lgxcdI{1} balance
five \lgxcdI{1}. Remove two \lgxcdI{1} from each pan: \lgxcdI{2} $= 3$ \lgxcdI{1}.
\step{3}{Balance (3): find \lgxcdI{4}}
\lgxcdI{4}\,\lgxcdI{1} balance two \lgxcdI{2}, i.e.\ $2\times 3 = 6$ \lgxcdI{1}. Remove one \lgxcdI{1} from each pan:
\lgxcdI{4} $= 5$ \lgxcdI{1}.
\step{4}{Part b)}
\lgxcdI{4}\,\lgxcdI{4}\,\lgxcdI{2}\,\lgxcdI{2} weigh as much as $5 + 5 + 3 + 3 = 16$ \lgxcdI{1}. Each \lgxcdI{3} equals $2$
\lgxcdI{1}, so we need $16 \div 2 = 8$ objects \lgxcdI{3}.
\step{5}{Part c)}
\lgxcdI{4}\,\lgxcdI{3}\,\lgxcdI{2} weigh as much as $5 + 2 + 3 = 10$ \lgxcdI{1}, i.e.\ $10\times 12 = 120$ g.
\begin{tipbox}If we call the masses $t$ (triangle), $q$ (square), $s$ (star), $o$ (circle), the three balances are three equations:
$q + t = 5o$; \ $2t = 4o$; \ $s + o = 2q$. Solving in turn gives $t = 2o$, $q = 3o$, $s = 5o$.\end{tipbox}
\verify{Balance (1): $3 + 2 = 5$. Balance (2): $2 + 2 = 4$. Balance (3): $5 + 1 = 6 = 3 + 3$. All three balances are balanced.}
\ans{a) \lgxcdI{3} $= 2$, \lgxcdI{2} $= 3$, \lgxcdI{4} $= 5$ objects \lgxcdI{1};\quad b) $8$ objects \lgxcdI{3};\quad c) $120$ g}
}

% ---------------------------------------------------------------- Problem 39
\begin{probbox}{39}{\relax}
The ``wheel'' below has $7$ circles: one at the hub and $6$ on the rim, forming $3$ straight ``spokes'' of $3$ circles each (two opposite
circles on the rim and the hub). Place the numbers $1, 2, 3, 4, 5, 6, 7$, one in each circle, so that the three numbers on each spoke
have the same sum $S$.
\begin{figC}
\foreach \a in {90,150,210}{\draw[Ink,line width=0.9pt] (\a:1.5)--(\a+180:1.5);}
\draw[Line,line width=0.8pt] (0,0) circle (1.5);
\foreach \a in {30,90,...,330}{\draw[Ink,fill=white,line width=0.9pt] (\a:1.5) circle (0.3);}
\draw[Ink,fill=tmBlue!12,line width=0.9pt] (0,0) circle (0.3);
\end{figC}
\begin{enumerate}[label=\alph*),leftmargin=1.6em,itemsep=1pt,topsep=2pt]
  \item Let $c$ be the number at the hub. Prove that $3S = 28 + 2c$.
  \item Deduce that the hub can only be $1$, $4$ or $7$. In each case find $S$ and give one arrangement.
\end{enumerate}
\end{probbox}
\probsol{
\textbf{Understanding the problem.} This is the ``count twice'' technique: add all the spokes and see how many times each circle is
counted.
\step{1}{Part a): add the three spokes}
The three spokes total $S + S + S = 3S$. In this sum every rim circle is counted \emph{once} (each rim circle lies on only one spoke),
while the hub is counted \emph{three} times. So
\[ 3S = (1 + 2 + \dots + 7) + 2c = 28 + 2c. \]
\begin{recallbox}$1 + 2 + 3 + 4 + 5 + 6 + 7 = 28$. The hub already appears once in the total $28$, so we only add $2c$ (the other two
times).\end{recallbox}
\step{2}{Part b): which values can $c$ take?}
$3S$ is divisible by $3$, so $28 + 2c$ must be divisible by $3$. Try $c = 1, 2, \dots, 7$:
\begin{center}\small
\begin{tabular}{|c|c|c|c|c|c|c|c|}
\hline
\rowcolor{tmBlue!10} $c$ & $1$ & $2$ & $3$ & $4$ & $5$ & $6$ & $7$\\ \hline
$28 + 2c$ & $\mathbf{30}$ & $32$ & $34$ & $\mathbf{36}$ & $38$ & $40$ & $\mathbf{42}$\\ \hline
Divisible by $3$? & yes & no & no & yes & no & no & yes\\ \hline
\end{tabular}
\end{center}
So $c \in \{1; 4; 7\}$, giving $S = 10$; $12$; $14$ respectively.
\step{3}{Concrete arrangements}
On each spoke the two opposite circles must add up to $S - c$, so we split the remaining $6$ numbers into $3$ pairs with equal sums.
\begin{itemize}[leftmargin=1.4em,itemsep=0pt,topsep=1pt]
  \item $c = 1$, $S = 10$: pairs with sum $9$: $(2, 7), (3, 6), (4, 5)$.
  \item $c = 4$, $S = 12$: pairs with sum $8$: $(1, 7), (2, 6), (3, 5)$.
  \item $c = 7$, $S = 14$: pairs with sum $7$: $(1, 6), (2, 5), (3, 4)$.
\end{itemize}
\begin{figC}
\foreach \cx/\cc/\p in {0/1/{2,3,4,7,6,5},4/4/{1,2,3,7,6,5},8/7/{1,2,3,6,5,4}}{
 \begin{scope}[xshift=\cx cm,scale=0.7]
  \foreach \a in {90,150,210}{\draw[Ink,line width=0.8pt] (\a:1.5)--(\a+180:1.5);}
  \foreach \v [count=\k from 0] in \p{\draw[Ink,fill=white,line width=0.8pt] (90+\k*60:1.5) circle (0.33);
     \node[font=\sffamily\small\bfseries,text=tmBlue] at (90+\k*60:1.5){\v};}
  \draw[Ink,fill=tmBlue!12,line width=0.8pt] (0,0) circle (0.33); \node[font=\sffamily\small\bfseries] at (0,0){\cc};
 \end{scope}}
\node[font=\small] at (0,-1.55){$S = 10$}; \node[font=\small] at (4,-1.55){$S = 12$}; \node[font=\small] at (8,-1.55){$S = 14$};
\end{figC}
\begin{warnbox}Do not assume ``the hub must be the middle number of the list ($4$)''. Both $1$ and $7$ can also go at the hub!\end{warnbox}
\verify{$S = 10$: $1 + 2 + 7 = 1 + 3 + 6 = 1 + 4 + 5 = 10$. \ $S = 12$: $4 + 1 + 7 = 4 + 2 + 6 = 4 + 3 + 5 = 12$. \
$S = 14$: $7 + 1 + 6 = 7 + 2 + 5 = 7 + 3 + 4 = 14$.}
\ans{a) see the proof;\quad b) $c = 1, S = 10$;\ \ $c = 4, S = 12$;\ \ $c = 7, S = 14$ (arrangements as in the figure)}
}

% ---------------------------------------------------------------- Problem 40
\begin{probbox}{40}{\stars{1}}
Place the numbers $1, 2, 3, 4, 5, 6$, one in each circle on the sides of the triangle (three circles at the vertices, three at the
midpoints of the sides), so that the three numbers on each side have the same sum $S$.
\begin{figC}
\draw[Ink,line width=0.9pt] (0,2.2)--(-2.2,-1.3)--(2.2,-1.3)--cycle;
\foreach \p in {(0,2.2),(-2.2,-1.3),(2.2,-1.3),(-1.1,0.45),(1.1,0.45),(0,-1.3)}{\draw[Ink,fill=white,line width=0.9pt] \p circle (0.3);}
\end{figC}
\begin{enumerate}[label=\alph*),leftmargin=1.6em,itemsep=1pt,topsep=2pt]
  \item Let $V$ be the sum of the three numbers at the vertices. Prove that $3S = 21 + V$. Hence find all possible values of $S$.
  \item For each value of $S$, give one arrangement.
  \item When $S = 10$, which three numbers are at the vertices? Explain why there is no other possibility.
\end{enumerate}
\end{probbox}
\probsol{
\textbf{Understanding the problem.} As in Problem 39, we ``count twice'': each vertex circle lies on \emph{two} sides, each midpoint
circle on \emph{one} side.
\step{1}{Part a): add the three sides}
The three sides total $3S$. In this sum each midpoint number is counted once and each vertex number twice. So
\[ 3S = (1 + 2 + 3 + 4 + 5 + 6) + V = 21 + V. \]
\step{2}{Bounds on $V$}
$V$ is the sum of three different numbers from $1, \dots, 6$: at least $1 + 2 + 3 = 6$, at most $4 + 5 + 6 = 15$. Also $21 + V$ is
divisible by $3$, and $21$ is divisible by $3$, so $V$ is divisible by $3$:
\[ V \in \{6;\ 9;\ 12;\ 15\} \quad\Rightarrow\quad S = \frac{21 + V}{3} \in \{9;\ 10;\ 11;\ 12\}. \]
\begin{recallbox}If $a$ is divisible by $3$ and $a + b$ is divisible by $3$, then $b$ is also divisible by $3$.\end{recallbox}
\step{3}{Part b): one arrangement for each $S$}
Choose three vertex numbers with sum $V$; then subtract the sum of the two end numbers of each side from $S$ to get the number at the
middle of that side.
\begin{figC}
\foreach \cx/\S/\t/\l/\r/\ml/\mr/\mb in {0/9/1/2/3/6/5/4, 3.6/10/1/3/5/6/4/2, 7.2/11/2/4/6/5/3/1, 10.8/12/4/5/6/3/2/1}{
 \begin{scope}[xshift=\cx cm,scale=0.62]
  \draw[Ink,line width=0.8pt] (0,2.2)--(-2.2,-1.3)--(2.2,-1.3)--cycle;
  \foreach \p/\v/\col in {(0,2.2)/\t/Ink,(-2.2,-1.3)/\l/Ink,(2.2,-1.3)/\r/Ink,(-1.1,0.45)/\ml/tmBlue,(1.1,0.45)/\mr/tmBlue,(0,-1.3)/\mb/tmBlue}{
    \draw[Ink,fill=white,line width=0.8pt] \p circle (0.36); \node[font=\sffamily\small\bfseries,text=\col] at \p {\v};}
  \node[font=\small] at (0,-2.2){$S = \S$};
 \end{scope}}
\end{figC}
Example for $S = 9$: vertices $1, 2, 3$; side $1$--$2$ needs $9 - 3 = 6$; side $1$--$3$ needs $9 - 4 = 5$; side $2$--$3$ needs
$9 - 5 = 4$.
\step{4}{Part c): $S = 10$ gives $V = 3\cdot 10 - 21 = 9$}
The triples of different numbers with sum $9$ are $\{1, 2, 6\}$, $\{1, 3, 5\}$, $\{2, 3, 4\}$. Try each one:
\begin{itemize}[leftmargin=1.4em,itemsep=0pt,topsep=1pt]
  \item Vertices $\{1, 2, 6\}$: the side joining $1$ and $2$ needs a middle number $10 - 3 = 7$ -- there is no $7$. \textbf{Ruled out.}
  \item Vertices $\{2, 3, 4\}$: the side joining $3$ and $4$ needs $10 - 7 = 3$ -- but $3$ is already at a vertex. \textbf{Ruled out.}
  \item Vertices $\{1, 3, 5\}$: the middle numbers are $10 - 4 = 6$, $10 - 6 = 4$, $10 - 8 = 2$ -- exactly the remaining numbers
  $2, 4, 6$. \textbf{Works.}
\end{itemize}
So when $S = 10$, the vertices must be $1, 3, 5$.
\begin{warnbox}The condition ``$V$ is divisible by $3$'' is only a \emph{necessary} condition: it tells us what $S$ \emph{could} be, but
we still have to show an actual arrangement (part b) and check each triple (part c).\end{warnbox}
\begin{parentbox}Let your child cut out $6$ slips of paper numbered $1$ to $6$ and try placing them on the drawing. After a few tries,
ask: ``How many times is a vertex number added?'' -- that is the main idea of the solution.\end{parentbox}
\verify{$S = 10$: $1 + 6 + 3 = 3 + 4 + 5 = 5 + 2 + 1 = 10$. $S = 12$: $4 + 3 + 5 = 5 + 1 + 6 = 6 + 2 + 4 = 12$.}
\ans{a) $S \in \{9; 10; 11; 12\}$;\quad b) as in the figure;\quad c) the vertices are $1, 3, 5$}
}
